\documentclass[11pt]{article}
\usepackage[margin=25mm]{geometry}
\usepackage{fontspec,unicode-math}
\setmainfont{Latin Modern Roman}
\setmathfont{Latin Modern Math}
\usepackage{amsmath,amssymb,mathtools}
\usepackage{xurl,hyperref}
\hypersetup{colorlinks=true,urlcolor=blue,linkcolor=blue}
\setcounter{secnumdepth}{0}
\setlength{\parindent}{0pt}
\setlength{\parskip}{5pt}
\setlength{\emergencystretch}{3em}
\newcommand{\sref}[2]{\hyperlink{source:#1:#2}{[S#2]}}
\newcommand{\cataloguescope}[1]{\paragraph{Source question.}#1}
\newcommand{\floor}[1]{\left\lfloor#1\right\rfloor}
\begin{document}
\section{10 Lectures and 42 Open Problems — Consistency and sample complexity of Multireference Alignment}\label{statement}

\subsection{Definitions and mathematical statement}

\paragraph{Shared notation.}Standard number systems are denoted by $\mathbb N,\mathbb Z,\mathbb Q,\mathbb R,\mathbb C$; index ranges, computational models and specialized notation follow the cited problem source and the definitions below.


Fix a signal length $L$ and unknown $u\in\mathbb{R}^L$. A cyclic shift acts by $(R_\ell u)_j=u_{j-\ell\bmod L}$. In the standard uniform-shift Gaussian model, independent observations are
\[
 y_i=R_{\ell_i}u+\sigma\xi_i,\qquad
 \ell_i\sim\operatorname{Unif}(\mathbb{Z}/L\mathbb{Z}),\quad
 \xi_i\sim N(0,I_L),
\]
with independent shifts and noise, and noise level $\sigma>0$. The signal is identifiable only up to a common cyclic shift; estimation error is therefore $\min_\ell\|\widehat u-R_\ell u\|_2$.

For the unitary discrete Fourier transform $\widehat u_k=L^{-1/2}\sum_j u_j e^{-2\pi i jk/L}$, the power spectrum is $(|\widehat u_k|^2)_k$; shifts leave it invariant. The noise contribution to the empirical power spectrum must be removed or otherwise fixed before judging phase consistency. Consistency means that the orbit error tends to zero in probability as the number of observations tends to infinity at fixed $L$ and $\sigma$. Sample complexity asks how many observations guarantee a specified error and success probability, as a function of signal class, $L$ and $\sigma$.

Two distinct likelihood conventions matter. A joint/profile estimator optimizes over unknown shifts as incidental parameters, minimizing $\sum_i\min_\ell\|y_i-R_\ell u\|_2^2$. A marginalized maximum-likelihood estimator instead maximizes
\[
 \prod_i\left(\frac1L\sum_{\ell=0}^{L-1}
 \exp\left[-\frac{\|y_i-R_\ell u\|_2^2}{2\sigma^2}\right]\right).
\]
These objectives need not have the same statistical behavior; the source lecture discussion refers to the quasi-MLE as well as the MLE.
\paragraph{Statement recorded in the supplied source.}
Is the Maximum likelihood for Multireference Alignment consistent? (after fixing the power spectrum) What is the sample complexity of the Multireference Alignment problem? \sref{statement}{2}
\paragraph{Scope qualification.}
This retains both parts of the historical question. The later primary sample-complexity paper establishes the cubic inverse-SNR scaling for generic signals in the high-noise uniform-shift regime; that is substantial progress on the second part and must not be described as still wholly unknown. Consistency of a marginalized likelihood and of a profile/quasi-likelihood are separate claims; the catalogue's short wording does not settle which estimator or signal class is intended. The author's original question and later primary literature were inspected on 2026-10-11, without claiming a complete resolution-status audit for every estimator or nongeneric case.
\paragraph{Source status.}
The supplied catalogue labels this record partially solved. Its source-status metadata was checked against the supplied snapshot on 2026-10-11. This import does not independently establish current open status or mathematical correctness.
\subsection{Short English statement}
After correcting the signal's Fourier magnitudes, does likelihood-based alignment consistently recover its phases, and how many noisy shifted samples are needed for accurate recovery?
\subsection{Sources}
\begin{itemize}
\item[\hypertarget{source:statement:1}{S1}] Ulam AI and the UnsolvedMath Contributors, UnsolvedMath: A Curated Collection of Open Mathematics Problems, catalogue entry. Catalogue curation and metadata: CC BY 4.0. The supplied source question is retained; definitions and the English summary are editorial additions.
\url{https://huggingface.co/datasets/ulamai/UnsolvedMath}.
\item[\hypertarget{source:statement:2}{S2}] AMR Open Problem Lists, original problem source linked by the catalogue
\url{https://afonsobandeira.wordpress.com/2015/12/09/42l10psynchronization/}.
\item[\hypertarget{source:statement:3}{S3}] Afonso S. Bandeira, Ten Lectures and Forty-Two Open Problems in the Mathematics of Data Science (2016), Section 10.3, Open Problem 10.4; Gaussian shift model and quasi-MLE qualification.
\url{https://people.math.ethz.ch/~abandeira/TenLecturesFortyTwoProblems.pdf#page=139}.
\item[\hypertarget{source:statement:4}{S4}] Amelia Perry, Jonathan Weed, Afonso S. Bandeira, Philippe Rigollet and Amit Singer, The Sample Complexity of Multireference Alignment (2017/2019); generic-signal sample complexity.
\url{https://arxiv.org/abs/1707.00943}.
\end{itemize}
\paragraph{Preparation and review.}
Original question attributed to the cited source and UnsolvedMath contributors. Codex (GPT-6; exact serving revision not exposed) assisted with definitions, English summaries, source inspection and formatting. No new human mathematical authorship or independent proof verification is asserted. Catalogue identity checked against the complete supplied registry and existing repository statements; known flagged repetitions and obvious identical conjectures excluded. Full mathematical-equivalence/literature review has not been performed. Statement preparation and final review followed contribution prompts C and D; automated checks verify structure and identity, not open status.
\label{end:statement}
\end{document}
